%% kapa: exam=1 round=31 year=2020 period=2 form=A subject=accounting subject_ko=회계학 seq=5 number=45
%% source: Q-Net 감정평가사 기출문제 (cst00309) artlSeq=5207219 file=extracted/1008/5207219/A형/A형/2020년도 제31회 감정평가사 1차 시험_2교시_A형.pdf
%% conversion: kordoc 4.24.51 -> tools/gen1.py (auto); review=proofread-sonnet; flags=none
\begin{question}{45}
(주)감평은 20x1년 초 임대수익을 얻고자 건물(취득원가 ￦1,000,000, 내용연수 5년, 잔존가치 ￦100,000, 정액법 상각)을 취득하고, 이를 투자부동산으로 분류하였다. 한편, 부동산 경기의 불황으로 20x1년 말 동 건물의 공정가치는 ￦800,000으로 하락하였다. 동 건물에 대하여 공정가치모형을 적용할 경우에 비해 원가모형을 적용할 경우 (주)감평의 20x1년도 당기순이익은 얼마나 증가 혹은 감소하는가? (단, 동 건물은 투자부동산의 분류요건을 충족하며, (주)감평은 동 건물을 향후 5년 이내 매각할 생각이 없다.)
\begin{choices}
  \item ￦20,000 증가
  \item ￦20,000 감소
  \item ￦0
  \item ￦180,000 증가
  \item ￦180,000 감소
\end{choices}
\end{question}
\answer{1}
